\section{Binary edge collisions in the narrow window}
\label{sec:binary-edge-poisson}

We now establish the analogous Poisson limit law for edge collisions in the binary hypercube. Let \(q=2\), \(N=n-1\), and \(s=N^{1/3}\). Evaluate the constants of Section~\ref{sec:binary-orbit-poisson} at \(q=2\), and set
\[
 u_{\rm edge}=\frac{3}{2R},\qquad
 K=\kappa_1\sqrt N\exp\left(As+\frac{t_2}{s}
                    +\frac{u_{\rm edge}\log N+y}{s^2}\right).
\]
For fixed \(y\), the number of colliding unordered pairs of distinct edges converges in distribution to twice a Poisson random variable of mean \(\Lambda_2(y)=\sqrt{2\pi v}\,e^{-D-Ry}/16\), with \(v=1/4\). In particular, the edge resolution probability tends to \(e^{-\Lambda_2(y)}\). This holds both under independent landmark selection with probability \(K/2^n\) and for a uniform random subset of size \(K\) rounded to an integer.

The argument reduces edge collisions to vertex collisions in dimension \(N\) by projecting each edge along its axis. Parity eliminates all nonparallel collisions and all odd-separation parallel collisions. On a single axis, the even-distance atom follows from Section~\ref{sec:binary-orbit-poisson}, and for tuples spanning multiple axes we project onto the common core of unshared coordinates to carry out the spanning-forest and Fourier factorization arguments.

\step{Parity and single atoms.}
We first show that nonparallel collisions and odd-separation parallel collisions occur with vanishing probability, and then compute the single-pair atom for parallel edges at a typical even distance.

Represent an edge by its axis \(i\) and its projected center in the \(N\)-cube obtained by deleting coordinate \(i\). Up to a sign depending on its center, the alternating sum of its shell counts equals the sum over all selected landmark words of \((-1)^{\sum_{h\ne i}z_h}\). Consequently, a collision between two nonparallel edges forces two such signed sums on distinct axes to agree up to sign. For either choice of sign, half of the cube contributes zero to the difference, and each nonzero sign is represented by a quarter of the cube. Under independent selection, the maximal binomial atom bound gives probability \(O(K^{-1/2})\) per comparison. In the uniform subset model, we use the iid draw model conditioned on no repetitions: condition first on the number of nonzero contributions, which is distributed as \(\mathrm{Bin}(K,1/2)\), and then apply the maximal atom bound to the independent fair signs. Since the probability of no repetition is \(1-o(1)\), the same bound holds, and a union bound over all pairs of axes and signs gives \(o(1)\). By the same argument, an odd-separation parallel collision forces the signed sum on its shared axis to equal zero, which occurs with probability \(O(K^{-1/2})\) per axis in either model.

On a single axis, each projected word has two preimages in the full cube, so in the Poisson occupancy model their counts simply add and the projected shell intensities are exactly \(K2^{-N}\binom Nr\). The single even-distance typical atom of Section~\ref{sec:binary-orbit-poisson} therefore applies with \(n\) replaced by \(N\). Compared to the vertex scaling, the new centering adds \(\log N\) to the scalar \(S\). This follows from the expansion in Section~\ref{sec:binary-orbit-poisson}, since the derivative in the displayed shift parameter is \(R+o(1)\), with an error that remains \(o(1)\) after multiplication by \(\log N\) by the clipped-sum and boundary expansions; equivalently, one can substitute the new coefficient of \(\log N\) directly into the scalar expansion. The resulting collision atom is
\[
 (1+o(1))\,N^{-1}2^{-2N}\frac{\sqrt{2\pi v}}2e^{-D-Ry}.
\]
The one-anchor bound from Section~\ref{sec:binary-orbit-poisson} removes all other even nonantipodal distances outside the typical distance band, even after multiplying by the \(n\) possible axes and, in the fixed-size model, by the conditioning penalty \(O(\sqrt K)\). Even projected antipodes have an expected count per axis bounded by \(\exp(-As/4+O(\log N))\), by the folded-shell atom of Section~\ref{sec:binary-orbit-poisson}. Its fixed-size relative transfer uses iid draws on the full cube: projection has the same one-draw discrepancy law, and excluding repetitions on full words changes the probability by only a relative error of \(o(1)\) by the shifted-atom estimate. Odd projected antipodes have already been excluded by parity.

\step{Cluster bounds across axes.}
Next we bound the contribution of clustered or atypical tuples of parallel edge pairs by restricting to the common core of coordinates outside their active axes.

Simultaneously complementing the two projected centers of a parallel edge pair reverses both histograms and leaves the collision event unchanged. We therefore count retained parallel collisions on these complement orbits, keeping the axis as part of the orbit label. For a fixed factorial tuple of order \(k\), let \(I\) be its set of active axes and put \(M=n-|I|\). Once the landmark letters on \(I\) are frozen, the distance to every edge in the tuple equals its common-core distance in dimension \(M\) plus a fixed offset in \([0,|I|-1]\). We cluster the endpoints on the common core using the same quotient-distance rule as in Section~\ref{sec:binary-orbit-poisson}. Choosing the active axes and the frozen endpoint letters costs only a polynomial factor. For \(b\) clusters in \(c\) connected components, the same entropy bound applies in dimension \(M\), and the contracted multigraph has no loops by the typicality of internal distances.

\step{Exact columns on the common core.}
To construct the exact columns needed for a rooted spanning forest, first align each close comparator to the anchor by complementing it when needed; absolute layers are measured from \(N/2\) and are unaffected by this operation. For the upper radius of a layer, choose all frozen letters seen by the anchor to mismatch it, giving the maximal offset; for the lower radius, choose them to match, giving offset zero. Every comparator offset is therefore weakly inward relative to the anchor's offset. If a nominated comparator uses a different axis, we can choose the anchor's free letter to match it at the upper radius or mismatch it at the lower radius, making its offset strictly inward. For a same-axis comparator with different aligned frozen letters, the offset is automatically strictly inward.

On the common core, tilt the mismatch probability to the requested core radius divided by \(M\), round the mismatch counts upward at the upper radius and downward at the lower radius in each tuple cell, and restore the total in a linear-size cell where all close aligned comparators agree with the anchor; such a cell exists once the fixed cluster parameter is taken sufficiently small. On the differing set of each close comparison, the mismatch count then lies on the inward side of half the set size. If a nominated comparison differs from the anchor on the core and still achieves equality, adjust one count on its differing set in the inward direction and compensate again in the common agreement cell. This preserves all weak inequalities.

Combining the core and offset comparisons gives weak inwardness for all close comparators and strict inwardness against any nominated distinct aligned edge. After excluding a sufficiently large fixed central strip, the contraction of the tilted mean together with bounded rounding and offsets prevents overshoot, and far comparisons are strictly inward by the same contraction bound. Applying Stirling's approximation to finitely many cells (including bounded-size cells and the fixed frozen pattern) gives polynomial relative shell mass. Consequently, the same leading-layer triangular argument and maximal binomial atom bound yield one complete family cost of \(2LN+O_k(N^{2/3}\log N)\) per forest link.

\step{Entropy comparison.}
For \(b>2c\), these \(b-c\) complete families defeat the configuration entropy bound. When \(b=2c\), every component has two clusters. If such a component has two candidate links whose aligned edges on some side are distinct, we anchor both links on that side and nominate the other anchor for strict inwardness. This works even for identical common-core words on different axes, and it supplies two complete families. If neither side has distinct aligned edges, the two links use the same axis and the same projected quotient endpoints. Distinct complement orbits then have opposite relative complement signs, and the third column per layer from the folded histogram constraint in Section~\ref{sec:binary-orbit-poisson} provides the extra half-family. In both cases the total cost defeats the entropy bound.

For the remaining single-link components, the spanning-forest bound also absorbs the \(\exp(-\Omega_k(N^{5/6}))\) fraction of common-core tuples whose coordinate-type frequencies differ from uniform by more than \(O_k(N^{11/12})\). All of these negligibility estimates absorb the fixed-size conditioning penalty.

\step{Relative typical joint atoms and counting.}
Finally, we show that the joint collision atoms of typical tuples factorize relatively into single-pair atoms, and we deduce the limiting Poisson distribution from the resulting factorial moments.

For typical tuples, carry out the Fourier singleton construction conditioned on a fixed frozen landmark pattern on \(I\). For any even projected pair, the difference between its frozen offsets is congruent modulo two to the distance between its common-core endpoints. Thus the full anchor radius and the deleted partner radius of matching parity impose compatible core radius constraints, and equality for every other pair likewise imposes compatible core difference constraints. The unimodular minor on linear-size common-core tuple cells is unchanged. Because the fixed offset shifts are bounded, the previous singleton intensity bound continues to hold, and omitting the anchor-radius constraint again gives a polynomial-mass neutral column. Moreover, conditioning on the frozen pattern preserves the two-tilt estimate for any two distinct typical endpoints: removing finitely many coordinates preserves the type and distance estimates and shifts radii only by bounded amounts. Thus the small-row splitting and dense-box estimates remain unchanged.

In the dense Gaussian comparison, the radial Walsh inner products vanish across different degrees. On the same axis, they are identical to the previous degree-wise inner products in dimension \(N\). Across different axes, retaining only equal Walsh supports from the two radial sums shows that the degree-\(j\) inner product equals the binary radial inner product on the \(N-1\) shared coordinates multiplied by \(\binom{N-1}{j}/\binom Nj=(N-j)/N\). Each even-distance block omits degrees zero and \(N\). On the shared core, the reflected degree zero at \(j=N-1\) contributes only \(1/N\), and the remaining squared bounds sum to \(o(1)\) by the same matching and replacement estimates. The cross-blocks therefore have whitened Hilbert--Schmidt norm \(o(1)\). Just as in the vertex proof, the joint determinant has relative error \(o(1)\), and the Fourier integrals factorize relatively. The neutral and singleton classes also verify the hypotheses of the relative iid and no-repeat transfer of Section~\ref{sec:fixed-cardinality}.

For each ordered list of \(k\) axes, before discarding the negligible counting fraction of atypical tuples, the parity restrictions select an exact fraction \(2^{-k}\) of all independent labeled projected endpoints. Every retained ordered orbit tuple has \(4^k\) labeled representatives. Consequently, the factorial moments are asymptotic to the corresponding powers of \(n2^{2N}/8\) times the single atom, and that product tends to \(\Lambda_2(y)\) because \(n/N\to1\). Since the retained physical count is exactly twice the orbit count, the exclusions above transfer the Poisson limit law to the full count.
